\documentclass[a4paper,11pt]{article}
\usepackage[a4paper,margin=2cm]{geometry}
\usepackage[T1]{fontenc}
\usepackage[utf8]{inputenc}
\usepackage[ngerman,english]{babel}
\usepackage{lmodern}
\usepackage{microtype}
\usepackage[table]{xcolor}
\usepackage{parskip}
\usepackage{enumitem}
\usepackage[most]{tcolorbox}
\usepackage{titlesec}
\usepackage{xparse}
\usepackage[hidelinks]{hyperref}
\definecolor{HeaderBlueDark}{RGB}{70,110,170}
\definecolor{RowBlueLight}{RGB}{235,243,255}
\definecolor{RowBlueDark}{RGB}{220,233,250}
\definecolor{SoftGray}{RGB}{90,90,90}
\setlength{\parindent}{0pt}
\setlist[itemize]{leftmargin=1.4em,itemsep=0.35em,topsep=0.35em}
\linespread{1.06}
\titleformat{\section}[block]{\Large\bfseries\color{HeaderBlueDark}}{}{0pt}{}
\titlespacing{\section}{0pt}{1.3em}{0.8em}
\titleformat{\subsection}[block]{\large\bfseries\color{HeaderBlueDark}}{}{0pt}{}
\titlespacing{\subsection}{0pt}{1.0em}{0.6em}
\newtcolorbox{exbox}{enhanced,breakable,colback=RowBlueLight,colframe=RowBlueDark,boxrule=0.4pt,arc=1.5mm,left=1.2mm,right=1.2mm,top=1mm,bottom=1mm,title={Beispiele \textnormal{(Examples)}},fonttitle=\bfseries,colbacktitle=RowBlueDark,coltitle=black}
\NewDocumentEnvironment{grammarentry}{mm+b}{%
\begin{tcolorbox}[enhanced,breakable,colback=white,colframe=HeaderBlueDark,boxrule=0.6pt,arc=1.5mm,left=1.5mm,right=1.5mm,top=1mm,bottom=1mm,colbacktitle=HeaderBlueDark,coltitle=white,fonttitle=\bfseries,title={#1\hfill\normalfont\footnotesize #2}]
#3
\end{tcolorbox}}{}
\newcommand{\GermanText}[1]{\textbf{Deutsch: }#1\par}
\newcommand{\EnglishText}[1]{{\small\color{SoftGray}\textbf{English: }#1\par}}
\newcommand{\ExampleItem}[2]{\item \textbf{#1}\\{\small\color{SoftGray}#2}}
\title{auf ... hin\\ \large auf + accusative phrase + hin}
\date{}
\begin{document}
\maketitle
\tableofcontents
\newpage
\section{Einordnung und Wortart / Classification and part of speech}
\begin{grammarentry}{auf ... hin}{Grammatische Einordnung / grammatical classification}
\GermanText{Grammatischer Status: Zirkumposition in weiterer Grammatik.}
\EnglishText{Grammatical status: Circumposition in broader grammar.}
\end{grammarentry}
\section{Struktur, Stellung und Rektion / Structure, position and case}
\begin{grammarentry}{Klammerbau / Framing structure}{Kasus / case}
\GermanText{Muster: auf + Akkusativgruppe + hin. Rektion: Akkusativ.}
\EnglishText{Pattern: auf + accusative phrase + hin. Case: accusative.}
\end{grammarentry}
\section{Bedeutung / Meaning}
\begin{grammarentry}{Grundbedeutung / Core meaning}{Semantik / semantics}
\GermanText{beschreibt einen Anlass, eine Reaktion auf etwas, die Ausrichtung auf einen Zweck oder eine Prüfung hinsichtlich eines Merkmals.}
\EnglishText{expresses a trigger, a response to something, orientation toward a purpose, or investigation for a property.}
\end{grammarentry}
\section{Verwendung und Register / Usage and register}
\begin{grammarentry}{Gebrauch / Usage}{Typische Kontexte / typical contexts}
\GermanText{Häufig folgt auf auf ... hin eine Handlung als Reaktion: auf eine Beschwerde hin. Auch fachsprachlich: auf eine Krankheit hin untersuchen.}
\EnglishText{The expression often precedes an action that follows as a response: in response to a complaint. It is also used technically: to test for a disease.}
\end{grammarentry}
\section{Beispiele / Examples}
\begin{exbox}
\begin{itemize}
\ExampleItem{Auf ihren Vorschlag hin änderte die Gruppe den Plan.}{The group changed the plan following her suggestion.}
\ExampleItem{Auf eine Beschwerde hin wurde die Anlage geprüft.}{The system was inspected following a complaint.}
\ExampleItem{Das Wasser wird auf Schadstoffe hin untersucht.}{The water is tested for pollutants.}
\ExampleItem{Er spart auf den Ruhestand hin.}{He is saving with retirement in mind.}
\end{itemize}
\end{exbox}
\section{Abgrenzung / Comparison with related forms}
\begin{grammarentry}{Vergleich / Comparison}{Unterschiede / differences}
\GermanText{Aufgrund ihrer Bitte und auf ihre Bitte hin können ähnlich sein. Auf ... hin hebt häufig den Reiz und die nachfolgende Reaktion stärker hervor.}
\EnglishText{Aufgrund ihrer Bitte and auf ihre Bitte hin may be similar. The latter often emphasizes the trigger and subsequent reaction.}
\end{grammarentry}
\section{Typischer Fehler / Typical mistake}
\begin{grammarentry}{Kasus und Form / Case and form}{Falsch versus richtig / wrong versus correct}
\GermanText{Falsch: auf seinem Vorschlag hin. Richtig: auf seinen Vorschlag hin.}
\EnglishText{Incorrect: auf seinem Vorschlag hin. Correct: auf seinen Vorschlag hin. Auf ... hin governs the accusative, not the dative.}
\end{grammarentry}
\section{Kurze Übung / Short exercise}
\begin{grammarentry}{Ergänze die richtige Form / Complete the phrase}{Selbstkontrolle / self-check}
\GermanText{Ergänze die Lücke: Auf ... (sein Hinweis) hin wurde die Datei aktualisiert.}
\EnglishText{Fill in the blank: Auf ... (sein Hinweis) hin wurde die Datei aktualisiert.}
\end{grammarentry}
\subsection{Lösung / Answer}
\begin{grammarentry}{Lösung / Answer}{Kasus prüfen / check the case}
\GermanText{Auf seinen Hinweis hin wurde die Datei aktualisiert.}
\EnglishText{The file was updated following his suggestion.}
\end{grammarentry}
\section{Zusammenfassung / Summary}
\begin{grammarentry}{Merksatz / Key point}{Form und Bedeutung / form and meaning}
\GermanText{auf ... hin — Akkusativ. beschreibt einen Anlass, eine Reaktion auf etwas, die Ausrichtung auf einen Zweck oder eine Prüfung hinsichtlich eines Merkmals.}
\EnglishText{auf ... hin — accusative. expresses a trigger, a response to something, orientation toward a purpose, or investigation for a property.}
\end{grammarentry}
\section{Quellen / Sources}
\begin{grammarentry}{Fachliche Einordnung / Classification}{IDS und Duden / IDS and Duden}
\GermanText{Für die Abgrenzung zwischen eindeutigen Zirkumpositionen und ähnlichen Konstruktionen siehe IDS (grammis).}
\EnglishText{For the distinction between unambiguous circumpositions and similar structures, see IDS (grammis).}
\url{https://grammis.ids-mannheim.de/terms/view/506}
\url{https://grammis.ids-mannheim.de/systematische-grammatik/1422}
\end{grammarentry}
\end{document}
